% ============================================================================
% VERSIÓN CORREGIDA del capítulo 4 (cohesión, redacción y variedad léxica).
% El contenido técnico, las tablas, las figuras, las etiquetas y las citas
% se conservan idénticos a la versión original; solo cambia la prosa.
% ============================================================================

% El argumento opcional define cómo aparece el título en el índice
% (los capítulos van en mayúsculas; los subtítulos quedan como están).
\section[RESULTADOS, CONCLUSIONES Y RECOMENDACIONES]{Resultados, conclusiones y recomendaciones}

\subsection{Resultados}

Los resultados de este capítulo siguen el orden de los tres objetivos específicos del proyecto, de manera que cada evidencia pueda rastrearse hasta una decisión de diseño, una acción de implementación y un criterio de verificación concretos. Organizarlos de este modo facilita distinguir entre los productos construidos, los contratos técnicos que los sostienen y el comportamiento observado durante la validación final; por ello, cada bloque describe los entregables obtenidos, la manera en que se articulan con la arquitectura general y las métricas o evidencias con las que se comprobó su funcionamiento.

% ── Objetivo 1 ──────────────────────────────────────────────────────────────
\subsubsection{Análisis del protocolo AVATecH e integración con ELAN}

El primer objetivo específico planteaba analizar el mecanismo de extensión AVATecH de ELAN, combinando el estudio de la documentación oficial con la ingeniería inversa del código fuente, con el fin de definir los contratos de interfaz REST y el esquema de instalación de modelos que hacen posible integrar servicios externos de inferencia en el flujo de anotación. De su cumplimiento se desprenden dos resultados que se complementan entre sí. El primero corresponde a los contratos: la solicitud y la respuesta de inferencia que viajan por HTTP (ejemplificadas en el Anexo~II) y el esquema \textbf{manifest.json} que gobierna la instalación de modelos (Anexo~I). El segundo es su materialización en código: cinco clases principales y un conjunto de DTO implementados en el paquete \textbf{mpi.eudico.\allowbreak{}client.\allowbreak{}annotator.\allowbreak{}recognizer.\allowbreak{}ai}, organizados según el Código~\ref{lst:estructura_clases_elan}; el inventario completo del paquete consta en el Anexo~IV. La pieza central de esta integración es el reconocedor \textbf{AIOrchestrationRecognizer}, que implementa la interfaz \textbf{Recognizer} de ELAN y hace que la herramienta delegue el procesamiento de video al middleware sin tocar la base del software; gracias a ello, ELAN conserva intacto su ciclo de trabajo mientras el procesamiento especializado queda encapsulado en un componente externo, en línea con el propósito con el que fue concebido el protocolo AVATecH \cite{auer2014avatech}.

En la Figura~\ref{fig:panel_elan} se aprecia el panel de configuración \textbf{AIOrchestrationRecognizerPanel} tal como aparece dentro de ELAN cuando el usuario selecciona el reconocedor \textit{AI Orchestration Middleware Bridge}. Dicho panel traduce en controles visibles parámetros que hasta entonces solo podían enviarse desde clientes externos ---dirección del servidor, video a analizar, modelo de IA instalado, nivel de anotación destino y formato de etiqueta---, de modo que la configuración necesaria para lanzar una inferencia quedó incorporada al mismo entorno donde el investigador revisa y edita sus anotaciones, sin pasos intermedios fuera de la herramienta.

\begin{figure}[H]
    \centering
    \includegraphics[width=0.75\textwidth]{Diagramas/Panel configuracion del reconocer integrado en ELAN.png}
    \caption{Panel de configuración del reconocedor integrado en ELAN}
    \label{fig:panel_elan}
\end{figure}

Por su parte, la Figura~\ref{fig:dialogo_modelos} recoge el diálogo \textbf{AIModelManagerDialog}, accesible desde el botón \textit{Gestionar modelos\ldots} del panel. Allí se concentran las operaciones de instalación, activación y desactivación de modelos sin que el usuario abandone el entorno de anotación: cada botón del diálogo corresponde a una operación HTTP del middleware que el investigador ya no necesita conocer, con lo cual la gestión del ciclo de vida de los modelos, resuelta internamente a través de endpoints REST, quedó expuesta al usuario final por medio de controles gráficos coherentes con la interfaz de ELAN.

\begin{figure}[H]
    \centering
    \includegraphics[width=0.75\textwidth]{Diagramas/Dialogo de gestion de modelos integrados en ELAN.png}
    \caption{Diálogo de gestión de modelos integrado en ELAN}
    \label{fig:dialogo_modelos}
\end{figure}

En cuanto al cliente HTTP \textbf{AIOrchestration\allowbreak{}MiddlewareClient}, implementado sobre \textbf{java.net.http.\allowbreak{}HttpClient}, este trabaja exclusivamente con HTTP/1.1 para asegurar la compatibilidad con Uvicorn y evitar que se negocie automáticamente HTTP/2 frente a un servidor que no lo ofrece por defecto \cite{java11_httpclient,uvicorn_docs}. Su método \textbf{installModel} construye a mano el cuerpo \textbf{multipart/form-data} que transporta los paquetes ZIP, sin recurrir a dependencias externas, mientras que \textbf{buildErrorMessage} convierte los códigos de error estructurados del middleware en mensajes legibles dentro de la interfaz gráfica de ELAN. La integración respeta el contrato AVATecH en sus tres puntos esenciales: el reconocedor informa el avance con \textbf{RecognizerHost.\allowbreak{}setProgress()}, entrega los segmentos con \textbf{insertSegmentation\allowbreak{}AsTier()} y atiende la cancelación del usuario a través de \textbf{stop()}. El procedimiento para compilar estas clases con Maven, registrar el reconocedor mediante el mecanismo SPI de Java e incorporar el JAR resultante a una instalación de ELAN~7.1 se detalla paso a paso en el Anexo~IV, de manera que la integración puede reproducirse sobre el código fuente oficial de la herramienta.

% ── Objetivo 2 ──────────────────────────────────────────────────────────────
\subsubsection{Núcleo del orquestador: API REST, cola de jobs y gestión de contenedores Docker}

El segundo objetivo específico exigía codificar el núcleo del orquestador con Python y virtualización ligera mediante Docker, incorporando control de concurrencia y gestión del ciclo de vida de los contenedores para sostener la estabilidad del sistema durante inferencias simultáneas. En respuesta a ese planteamiento se desarrolló el servicio central del middleware, que concentra la validación de solicitudes, la exposición de endpoints, la coordinación de jobs y la entrega de métricas operativas. Optar por una API REST hizo que ELAN, Postman y cualquier otro cliente conversaran con el mismo contrato HTTP sin depender del lenguaje interno de cada componente, en consonancia con el principio de interfaces desacopladas propio de las arquitecturas basadas en recursos \cite{fielding2000rest,fastapi2026features}.

La API expone ocho endpoints organizados en cuatro grupos funcionales, resumidos en la Tabla~\ref{tab:endpoints_api}. El servidor Uvicorn sirve la aplicación FastAPI sobre el estándar ASGI y mantiene disponibles los endpoints de gestión ---entre ellos \textbf{GET /health} y \textbf{GET /api/v1/metrics}--- incluso mientras transcurre una inferencia prolongada. A ello se suma que el resultado de cada job permanece consultable durante la sesión a través de \textbf{GET /api/v1/jobs/\{job\_id\}}, gracias a un almacén en memoria (\textbf{MemoryStore}) que conserva las respuestas completadas hasta el siguiente reinicio del servicio. Este conjunto pone de manifiesto que el middleware no se limita a ejecutar inferencias: ofrece además las operaciones de observabilidad y administración necesarias para verificar el estado del servicio mientras un backend de IA consume recursos.

El componente \textbf{JobQueue} implementa una cola FIFO con límite configurable de concurrencia apoyándose en \textbf{collections.deque} y \textbf{threading.Event}, primitivas estándar de Python empleadas para coordinar espera y notificación entre hilos \cite{python_threading}. Cuántos backends Docker pueden inferir a la vez lo decide el parámetro \textbf{MIDDLEWARE\_\allowbreak{}MAX\_\allowbreak{}CONCURRENT\_\allowbreak{}JOBS}, ajustable desde el \textbf{docker-compose.yml} descrito en el Anexo~III. Con el valor por defecto de~1, un solo contenedor carga pesos en memoria en cada instante, lo que reduce el riesgo de desbordamiento bajo cargas concurrentes; entretanto, mientras un job aguarda su turno, el resto de la API continúa respondiendo. La Tabla~\ref{tab:metricas_concurrencia} recoge los resultados de la prueba controlada con dos solicitudes simultáneas y deja ver que la segunda permanece en cola hasta que la primera libera el slot de ejecución.

\begin{table}[H]
\centering
\caption{Métricas del sistema durante la ejecución de dos inferencias simultáneas}
\label{tab:metricas_concurrencia}
\small
\renewcommand{\arraystretch}{1.2}
\begin{tabular}{@{}>{\raggedright\arraybackslash}p{4.5cm}>{\centering\arraybackslash}p{2.5cm}>{\raggedright\arraybackslash}p{5.5cm}@{}}
\toprule
\textbf{Métrica / Parámetro} & \textbf{Valor real} & \textbf{Interpretación} \\
\midrule
\texttt{total\_jobs} & 10 & Total de solicitudes procesadas en el histórico. \\
\texttt{completed\_jobs} & 7 & Inferencias finalizadas con éxito. \\
\texttt{active\_jobs} & 1 & Inferencia en ejecución activa en el contenedor. \\
\texttt{queued\_jobs} & 1 & Segunda inferencia en espera dentro de la cola. \\
\texttt{timeout\_jobs} & 1 & Trabajos que excedieron el tiempo límite establecido. \\
\texttt{failed\_jobs} & 0 & Trabajos finalizados con error crítico del sistema. \\
\texttt{average\_exec\_ms} & 33673.43 & Tiempo promedio de ejecución global [ms]. \\
\texttt{last\_exec\_ms} & 32428.00 & Tiempo de respuesta de la última inferencia [ms]. \\
\texttt{DOCKER\_TIMEOUT} & 1 & Contador de errores por expiración del contenedor. \\
\bottomrule
\end{tabular}
\end{table}

Los valores de la Tabla~\ref{tab:metricas_concurrencia} respaldan cuantitativamente tanto el aislamiento como el ordenamiento de peticiones del middleware. Al llegar dos solicitudes de inferencia en el mismo instante, los contadores internos pasaron a \texttt{active\_jobs = 1} y \texttt{queued\_jobs = 1}: el mecanismo de exclusión mutua retuvo la segunda petición en la cola FIFO mientras el contenedor procesaba la primera, con lo cual se evitó saturar los recursos del equipo. Los contadores, además, resultan coherentes entre sí, pues los 10 jobs del histórico se reparten entre 7 completados, 1 expirado por timeout y los 2 de la prueba en curso (1 activo y 1 en cola). El tiempo registrado en \texttt{last\_exec\_ms}, cercano a 32.4 segundos, aporta a su vez una línea base del rendimiento del modelo de reconocimiento integrado.

La Figura~\ref{fig:postman_health} corresponde al escenario E-01 de la validación final y presenta la respuesta del endpoint \textbf{GET /health} verificada desde Postman. Comprobar esta disponibilidad básica antecede, en el orden de las pruebas, a las operaciones más costosas del sistema, como la instalación de modelos o las inferencias sobre video.

\begin{figure}[H]
    \centering
    \includegraphics[width=0.95\textwidth]{Diagramas/Respuesta del endpoint GET -health.png}
    \caption{Respuesta del endpoint \textbf{GET /health} verificada desde Postman (escenario E-01)}
    \label{fig:postman_health}
\end{figure}

En la Figura~\ref{fig:postman_metrics} puede observarse la respuesta del endpoint \textbf{GET /api/v1/metrics} después de ejecutar dos inferencias exitosas y un timeout forzado. La lectura de los contadores \textbf{total\_jobs}, \textbf{completed\_jobs} y \textbf{timeout\_jobs} deja constancia de que el middleware, además de procesar solicitudes, conserva trazabilidad cuantitativa del desenlace de cada job; esta información sirvió como evidencia de observabilidad durante la validación final.

\begin{figure}[H]
    \centering
    \includegraphics[width=0.95\textwidth]{Diagramas/Metricas acumuladas del middleware tras la prueba de carga controlada.png}
    \caption{Métricas acumuladas del middleware tras la prueba de carga controlada (\textbf{GET /api/v1/metrics})}
    \label{fig:postman_metrics}
\end{figure}

La gestión del ciclo de vida de los contenedores, segunda vertiente de este objetivo, tomó forma en el componente \textbf{DockerRunner} y en el contrato de comunicación entre el middleware y el backend del modelo. Este resultado enlaza la arquitectura de orquestación con el procesamiento real de video: el middleware no interpreta directamente las señas; prepara el entorno, invoca el backend de IA, valida su respuesta y adapta los segmentos al formato que ELAN puede insertar en la línea de tiempo.

En concreto, el \textbf{DockerRunner} traduce la ruta del video del sistema anfitrión al punto de montaje interno del contenedor, \textbf{/data/videos}, verifica la disponibilidad del backend mediante polling sobre \textbf{GET /health} con reintentos cada 250~ms y transforma la respuesta JSON del backend en objetos \textbf{TemporalSegment} validados por Pydantic \cite{pydantic2026models}. Conviene subrayar que el aporte del TIC en este punto no es el modelo de reconocimiento en sí, sino la capa que hace posible integrarlo: el middleware orquesta la ejecución sin conocer la arquitectura interna del pipeline, y cualquier backend que respete el contrato de \textbf{/health}, \textbf{/infer} y \textbf{manifest.json} puede ocupar ese lugar. Lo que sí corresponde al middleware es el postprocesamiento de las salidas, esto es, la validación y adaptación de la respuesta del backend en anotaciones semánticas compatibles con ELAN. Para la validación se empleó el backend \textbf{lsec\_bio\_gloss\_final\_v1}, cuyo desarrollo queda fuera del alcance de este trabajo; su pipeline interno corre encapsulado en un contenedor Docker, con las dependencias aisladas tanto de ELAN como del middleware \cite{docker2026overview}. El manifest completo con el que se instala este paquete se recoge en el Anexo~I.

La Figura~\ref{fig:postman_inferencia} ilustra la respuesta de una inferencia exitosa ejecutada desde Postman contra el endpoint \textbf{POST /api/v1/jobs/segment-video}, con segmentos detectados, marcas temporales, etiquetas de glosa, predicciones alternativas y trazabilidad de ejecución. Su relevancia radica en que vincula el contrato REST con el resultado lingüístico esperado, al constatar que el backend devuelve datos suficientes para que el middleware construya anotaciones temporales compatibles con ELAN. Un ejemplo completo de esta respuesta, con la descripción campo por campo del objeto \textbf{trace} y de sus etapas, figura en el Anexo~II.

\begin{figure}[H]
    \centering
    \includegraphics[width=0.95\textwidth]{Diagramas/Respuesta de inferencia exitosa.png}
    \caption{Respuesta de inferencia exitosa desde Postman: segmentos con glosas y trazabilidad completa}
    \label{fig:postman_inferencia}
\end{figure}

Los tiempos reales observados en el campo \textbf{trace.stages} de la respuesta del middleware durante la inferencia de validación se resumen en la Tabla~\ref{tab:latencia_etapas}. Los valores corresponden a un video de LSEc con una duración de 8.51~segundos (510 fotogramas a 59.94~FPS) procesado en un equipo con CPU Intel Core i5 de 12.ª generación. Desagregar la latencia por etapas ayuda a distinguir el costo de validación, espera en cola, arranque del contenedor, comprobación de disponibilidad, inferencia y postprocesamiento; de esa desagregación se desprende que la sobrecarga introducida por la orquestación es marginal, pues la suma de validación, cola y postprocesamiento apenas alcanza unos pocos milisegundos, mientras que el tiempo restante transcurre dentro del backend y depende del modelo instalado, no del middleware. El arranque del contenedor registró 0~ms porque este ya se encontraba activo (\textit{warm start}), condición habitual después de la primera ejecución.

\begin{table}[H]
\centering
\caption{Latencia por etapa de inferencia (modelo \textbf{lsec\_bio\_gloss\_final\_v1})}
\label{tab:latencia_etapas}
\small
\renewcommand{\arraystretch}{1.15}
\begin{tabular}{@{}>{\raggedright\arraybackslash}p{4.4cm}>{\centering\arraybackslash}p{3.0cm}>{\raggedright\arraybackslash}p{5.0cm}@{}}
\toprule
\textbf{Etapa} & \textbf{Latencia medida (ms)} & \textbf{Descripción} \\
\midrule
Validación del modelo          & 1        & Consulta al registro de modelos \\
Espera en cola (\texttt{QUEUED})   & 1        & Sin contención con 1 job activo \\
Arranque del contenedor         & 0        & Contenedor previamente activo (\textit{warm start}) \\
Comprobación de disponibilidad del backend        & 4        & Polling hasta disponibilidad \\
Inferencia en el backend        & 34\,181  & Procesamiento de 510 fotogramas \\
Postprocesamiento               & 1        & Validación y adaptación de respuesta \\
\textbf{Total (wall-clock)}    & \textbf{34\,198} & Desde \texttt{RECEIVED} hasta \texttt{COMPLETED} \\
\bottomrule
\end{tabular}
\end{table}

Como cierre del flujo, la Figura~\ref{fig:elan_anotaciones} presenta las anotaciones generadas por el sistema en la línea de tiempo de ELAN tras ejecutar el reconocedor desde el panel integrado, de acuerdo con los escenarios E-07 y E-08. Su valor probatorio está en que completa el ciclo: la inferencia ocurre fuera de ELAN, pero los resultados regresan como anotaciones temporales ubicadas en el nivel de anotación configurado. Las marcas de inicio y fin de cada anotación deben coincidir con los intervalos de señas visibles en el video procesado, ya que esa correspondencia confirma que la herramienta de anotación interpretó correctamente la segmentación temporal.

\begin{figure}[H]
    \centering
    \includegraphics[width=0.95\textwidth]{Diagramas/Anotaciones generadas automáticamente en ELAN por el reconocedor integrado.png}
    \caption{Anotaciones generadas automáticamente en ELAN por el reconocedor integrado (escenario E-07 y E-08)}
    \label{fig:elan_anotaciones}
\end{figure}

% ── Objetivo 3 ──────────────────────────────────────────────────────────────
\subsubsection{Validación técnica mediante pruebas funcionales}

El tercer objetivo específico, orientado a validar el desempeño técnico del middleware mediante pruebas de carga en escenarios de anotación de LSEc, se cubrió con quince escenarios distribuidos en cuatro categorías: conectividad, gestión de modelos, inferencia y manejo de errores. La validación midió latencia de comunicación, consumo de recursos computacionales y recuperación ante fallos, una combinación que abarcó tanto el comportamiento observable desde las interfaces visibles para el usuario como la consistencia interna de archivos, métricas, contenedores y rutas montadas.

La Tabla~\ref{tab:escenarios_validacion} sintetiza los escenarios más representativos de la validación final. Los identificadores E-01 a E-15 funcionan como etiquetas de trazabilidad que enlazan cada acción de prueba con su respuesta HTTP esperada, el resultado observado y la evidencia obtenida durante la ejecución.

\begin{table}[H]
\centering
\caption{Resultados de escenarios de validación final (selección representativa)}
\label{tab:escenarios_validacion}
\scriptsize
\renewcommand{\arraystretch}{1.15}
\setlength{\tabcolsep}{4pt}

\begin{tabular}{
    >{\centering\arraybackslash}p{1.0cm}
    >{\raggedright\arraybackslash}p{3.2cm}
    >{\centering\arraybackslash}p{1.5cm}
    >{\centering\arraybackslash}p{1.6cm}
    >{\raggedright\arraybackslash}p{5.4cm}
}
\toprule
\textbf{ID} & \textbf{Escenario} & \textbf{HTTP esp.} & \textbf{Resultado} & \textbf{Observación} \\
\midrule

E-01 &
Comprobación de disponibilidad desde Postman &
200 &
$\checkmark$ &
Respuesta menor a 50 ms. \\

E-02 &
Comprobación de disponibilidad desde ELAN &
200 &
$\checkmark$ &
El panel muestra el estado de servidor disponible. \\

E-03 &
Instalación de modelo mediante ZIP &
200 &
$\checkmark$ &
Imagen Docker construida y registry actualizado. \\

E-04 &
Listado de modelos &
200 &
$\checkmark$ &
Modelo visible con estado \texttt{available}. \\

E-05 &
Activar o desactivar modelo &
200 &
$\checkmark$ &
La solicitud \texttt{PATCH} actualiza el estado en el registry. \\

E-06 &
Instalación duplicada &
409 &
$\checkmark$ &
Error controlado \texttt{MODEL\_ALREADY\_EXISTS}. \\

E-07 &
Inferencia exitosa desde ELAN &
200 &
$\checkmark$ &
Anotaciones generadas en el nivel de anotación \texttt{AUTO\_GLOSS\_SEGMENTS}. \\

E-08 &
Verificación visual de anotaciones &
N/A &
$\checkmark$ &
Segmentos alineados temporalmente en la línea de tiempo. \\

E-09 &
Inferencia exitosa desde Postman &
200 &
$\checkmark$ &
Campo \texttt{trace} con tiempos por etapa. \\

E-10 &
Video inexistente &
502 &
$\checkmark$ &
Error controlado \texttt{DOCKER\_INFERENCE\_ERROR}. \\

E-11 &
Modelo desactivado &
409 &
$\checkmark$ &
Error controlado \texttt{MODEL\_DISABLED}. \\

E-12 &
Middleware detenido desde ELAN &
N/A &
$\checkmark$ &
ELAN informa que el servicio no está disponible. \\

E-13 &
Timeout forzado con \texttt{timeout\_sec = 1} &
504 &
$\checkmark$ &
Error controlado \texttt{DOCKER\_TIMEOUT}. \\

E-14 &
Contenedor del modelo detenido &
200 &
$\checkmark$ &
DockerRunner reinicia el contenedor automáticamente. \\

E-15 &
Métricas después de fallos &
200 &
$\checkmark$ &
Contadores \texttt{timeout\_jobs} y \texttt{error\_counts} actualizados correctamente. \\

\bottomrule
\end{tabular}
\end{table}

Por otro lado, la Tabla~\ref{tab:consumo_recursos} condensa el consumo de recursos medido con \textbf{docker stats} durante la inferencia completa con el modelo \textbf{lsec\_bio\_gloss\_final\_v1}, procesando el mismo video de 8.51 segundos. La medición se realizó a nivel de contenedores porque el middleware y el backend consumen recursos de manera muy distinta: el primero mantiene un perfil bajo ($\sim$0.13\,\% de CPU y $\sim$52.45\,MiB de RAM), asociado estrictamente a la orquestación, la validación y las métricas, mientras que el gasto intensivo se concentra en el contenedor del backend, que utiliza $\sim$194.89\,\% de CPU (casi dos núcleos completos) y $\sim$1.4\,GiB de memoria para ejecutar el pipeline de inferencia visual \cite{docker_cli_stats}. Semejante asimetría respalda una decisión de diseño central del sistema: dado que el costo computacional vive en el backend, es posible detener o reemplazar modelos sin comprometer la disponibilidad del orquestador.

\begin{table}[H]
\centering
\caption{Consumo de recursos durante inferencia (medición real con \texttt{docker stats})}
\label{tab:consumo_recursos}
\small
\renewcommand{\arraystretch}{1.15}
\begin{tabular}{@{}>{\raggedright\arraybackslash}p{3.8cm}>{\centering\arraybackslash}p{2.8cm}>{\centering\arraybackslash}p{2.8cm}>{\raggedright\arraybackslash}p{3.2cm}@{}}
\toprule
\textbf{Contenedor} & \textbf{CPU (\%)} & \textbf{RAM} & \textbf{Observación} \\
\midrule
Middleware (\texttt{elan-ai-middleware}) & 0.13 & 52.45\,MiB & Rol de orquestación \\
Backend del modelo (\texttt{elan-ai-model-...}) & 194.89 & 1.4\,GiB & Pipeline de inferencia \\
\bottomrule
\end{tabular}
\end{table}

En el plano interno, los escenarios de caja blanca corroboraron la consistencia del estado después de las operaciones principales: el archivo \textbf{registry.json} reflejó la entrada correcta del manifest tras la instalación; el mecanismo de bootstrap manifests reconstruyó el registro de forma automática después de eliminar \textbf{registry.json} y reiniciar el middleware; la traducción de la ruta del video quedó asentada en los logs del backend bajo el prefijo \textbf{/data/videos}; y \textbf{docker inspect} verificó el montaje correcto del directorio de videos. Estas comprobaciones complementan los códigos HTTP porque examinan que el estado persistente y la infraestructura Docker permanezcan coherentes tras cada operación.

Vistos en conjunto, los tres bloques de resultados recorren el ciclo completo del sistema: el primero prueba que ELAN puede delegar el procesamiento sin perder su flujo de anotación y que los contratos definidos sostienen esa comunicación; el segundo, que el orquestador regula la carga, administra los contenedores y devuelve la inferencia convertida en anotaciones temporales utilizables; y el tercero, que el comportamiento se mantiene predecible tanto en operación normal como frente a fallos. Los anexos I~a~IV acompañan estos resultados con los artefactos necesarios para reproducirlos: el contrato de empaquetado, un ejemplo íntegro de respuesta, la configuración de despliegue y el procedimiento de compilación e integración en ELAN.

\vspace{1\baselineskip}


\subsection{Conclusiones}

El trabajo realizado confirmó que es posible acercar servicios de Inteligencia Artificial al flujo de anotación de ELAN sin alterar en profundidad su arquitectura. La solución desarrollada opera como una capa intermedia entre la herramienta de anotación y los modelos de reconocimiento, de suerte que el investigador conserva su entorno habitual de trabajo mientras accede a resultados automáticos que después puede revisar y corregir. A partir de los objetivos planteados y de la validación efectuada se establecen las siguientes conclusiones.

\begin{enumerate}
    \item El análisis de ELAN y de su mecanismo de extensión permitió comprender cómo debe incorporarse un proceso externo a la línea de tiempo de anotación. Sin esa revisión, la solución habría terminado funcionando como una herramienta aislada; con ella, en cambio, se consiguió un apoyo integrado al proceso que el investigador ya realiza. Lejos de reemplazar el criterio humano, el sistema entrega segmentos y etiquetas preliminares que pueden aceptarse, ajustarse o descartarse dentro de ELAN, de manera que la integración respeta el uso lingüístico de la herramienta y deja la anotación final bajo responsabilidad del especialista.

    \item La arquitectura propuesta demostró que separar la anotación, la orquestación y la inferencia facilita conjugar tecnologías distintas sin comprometer la estabilidad del conjunto. ELAN se mantiene como espacio de revisión y edición; el middleware organiza las solicitudes, valida la información y coordina la ejecución; los modelos, por su parte, procesan los videos en un entorno independiente. Esta división reduce el acoplamiento entre componentes, impide que las dependencias de los modelos alcancen a ELAN y deja una base ordenada sobre la cual incorporar nuevos servicios de reconocimiento en el futuro.

    \item Apostar por una arquitectura local protegió los videos y anotaciones del investigador, puesto que la información no necesita salir de la estación de trabajo para ser procesada. A ello se añade que el control sobre la ejecución de inferencias evitó que varios procesos pesados compitieran sin coordinación por los mismos recursos del equipo, algo especialmente relevante en escenarios de investigación donde los equipos disponibles suelen tener capacidades limitadas y donde la estabilidad de la herramienta de anotación importa tanto como la precisión del modelo utilizado.

    \item Las pruebas realizadas mostraron que el middleware responde de manera controlada tanto en situaciones esperadas como frente a fallos comunes: videos inexistentes, modelos no disponibles, tiempos de espera excedidos o contenedores detenidos. La validación no se redujo a comprobar respuestas exitosas; hizo posible observar, además, cómo el sistema conserva su estado, informa los errores y mantiene abierta la consulta de progreso y métricas. Queda confirmado, con ello, que la solución no solo ejecuta inferencias, sino que ofrece condiciones mínimas de trazabilidad y diagnóstico para un uso real.

    \item La principal contribución del TIC es una arquitectura local, extensible y desacoplada que acerca los modelos de Inteligencia Artificial al trabajo cotidiano de anotación en Lengua de Señas Ecuatoriana. El sistema no pretende sustituir la revisión lingüística ni resolver por sí solo el problema del reconocimiento automático; ofrece, más bien, una vía práctica para que los resultados de esos modelos lleguen a ELAN como anotaciones temporales editables. Se abre así la posibilidad de reducir tareas repetitivas, apoyar la exploración de material audiovisual y facilitar investigaciones futuras que requieran integrar nuevos modelos sin reconstruir la herramienta completa.
\end{enumerate}

\vspace{1\baselineskip}

\subsection{Recomendaciones}

Los resultados obtenidos indican que el sistema puede seguir creciendo sin perder la idea central que guio este trabajo: mantener a ELAN como herramienta principal de anotación y usar el middleware como apoyo para integrar modelos externos de forma controlada. Las recomendaciones que siguen proponen mejoras futuras encaminadas a fortalecer la estabilidad, facilitar el mantenimiento y aumentar la utilidad de la solución en proyectos reales de documentación y análisis de LSEc.

\begin{enumerate}
    \item Convendría afinar la administración de los recursos de ejecución para que el sistema se adapte a equipos con capacidades diversas. Una línea de trabajo posible es incorporar un planificador que considere si cada modelo requiere CPU, GPU o una cantidad mayor de memoria antes de iniciar una inferencia; con ese criterio, los modelos livianos podrían ejecutarse con mayor flexibilidad, mientras los más exigentes seguirían protegidos por límites que eviten saturar el equipo del investigador.

    \item El diseño cliente-servidor del middleware abre una vía de crecimiento que merece aprovecharse: trasladar el orquestador y los modelos a un equipo de altas prestaciones, por ejemplo un servidor institucional con GPU, para que opere como servidor de inferencias compartido. Puesto que la comunicación entre ELAN y el middleware ocurre por HTTP, varias estaciones de trabajo de un mismo laboratorio podrían dirigir sus solicitudes a ese servidor apuntando la URL del reconocedor a la dirección IP correspondiente, sin modificar el código de ninguno de los dos lados. Dar ese salto exige ajustes de configuración y seguridad: publicar el puerto fuera de la interfaz local, ampliar el valor de \texttt{MIDDLEWARE\_MAX\_CONCURRENT\_JOBS} según la capacidad del equipo (con el valor actual de 1, las solicitudes simultáneas se atienden una por una en orden de llegada), centralizar el repositorio de videos y añadir autenticación y cifrado si la red no es de confianza. La misma idea admite extenderse a un despliegue en la nube; en tal caso, la decisión debería sopesar la soberanía de los datos que motivó la arquitectura local, porque los videos de LSEc saldrían de la infraestructura del investigador.

    \item Es aconsejable robustecer el almacenamiento local de la información de modelos instalados. Aunque el archivo actual cumple su función, una base de datos ligera registraría los cambios con mayor seguridad, facilitaría la consulta de versiones y reduciría el riesgo de inconsistencias ante una detención inesperada del servicio. Eso sí, la mejora debería mantenerse simple y local, de modo que no complique la instalación ni contradiga el principio de soberanía de los datos planteado en el proyecto.

    \item También se debería ampliar el conjunto de modelos probados con la arquitectura desarrollada, documentando y empaquetando nuevos modelos de reconocimiento de LSEc bajo el mismo formato de instalación usado en este TIC. Semejante continuidad permitiría comprobar la flexibilidad real del middleware frente a modelos con enfoques distintos, por ejemplo aquellos orientados a señas específicas, glosas más amplias o información no manual como expresiones faciales y postura corporal.

    \item Otra mejora valiosa apunta a la interacción del investigador con los resultados generados por la IA. Una extensión futura del panel en ELAN podría presentar, junto a la etiqueta principal, las alternativas de glosa cuando el modelo las entregue; la revisión ganaría así en transparencia y utilidad, porque el especialista contaría con más información para decidir si acepta una anotación, la corrige o la descarta antes de incorporarla definitivamente a la línea de tiempo.

    \item Finalmente, se recomienda complementar la validación técnica con una evaluación lingüística a cargo de especialistas en LSEc, que compare las anotaciones generadas automáticamente con anotaciones de referencia y considere no solo si la etiqueta propuesta es correcta, sino también si los tiempos de inicio y fin corresponden con el evento observado en el video. Sus resultados permitirían medir el aporte real del sistema en tareas de documentación lingüística y orientar mejoras tanto en los modelos como en la manera en que sus salidas se presentan al investigador.
\end{enumerate}
